% !TEX root = ../../main_docker.tex
% ---------------------------------------------------------------------
%  Referenz: alle Kommandos und Dockerfile-Instruktionen auf einen Blick,
%  plus die Rechnungen, die an der Pruefung drankommen koennen.
% ---------------------------------------------------------------------
\section{Referenz}
\label{sec:ref}

\subsection{Kommandos nach Aufgabe}
\begin{ztab}{@{}lX@{}}
\multicolumn{2}{@{}l}{\textbf{Images}}\\
\R{images}, \R{images -a} & lokale Images, \R{-a} auch namenlose \\
\R{pull <tag>} & Image holen, ohne zu starten \\
\R{build . -t <tag>} & Image aus Dockerfile im Kontext \R{.} \\
\R{build . -f <datei>} & anderes Dockerfile, Kontext bleibt \R{.} \\
\R{build . --target <stufe>} & nur bis zu dieser Build-Stufe \\
\R{history <image>} & Schichten mit Instruktion und Grösse \\
\R{inspect <image>} & \R{RootFS.Layers}, \R{Config.Cmd}, \R{Config.Entrypoint} \\
\R{tag <alt> <neu>} & zusätzlicher Name, Image bleibt dasselbe \\
\R{rmi <image>} & Image löschen \\
\R{image prune [-a]} & namenlose [bzw.\ alle unbenutzten] Images weg \\
\R{login -u <user>} / \R{image push} & auf Docker Hub anmelden / hochladen \\
\addlinespace
\multicolumn{2}{@{}l}{\textbf{Container}}\\
\R{run <image> [cmd]} & \emph{neuen} Container erzeugen und starten \\
\quad\R{-it} & interaktiv mit Terminal \\
\quad\R{-d} / \R{-itd} & im Hintergrund \\
\quad\R{--rm} & nach dem Stoppen automatisch löschen \\
\quad\R{--name <n>} & eigener Name statt Zufallsname \\
\quad\R{-p <host>:<cnt>} & Port durchreichen \\
\quad\R{-v <quelle>:<ziel>} & Volume einhängen \\
\R{ps}, \R{ps -a}, \R{ps -a -{}-size} & laufende / alle / mit Platzbedarf \\
\R{start}, \R{stop}, \R{restart} & vorhandenen Container steuern \\
\R{attach <cnt>} & sich an einen laufenden Container hängen \\
\R{exec -it <cnt> bash} & \emph{zusätzlichen} Prozess im Container \\
\R{logs [-f] <cnt>} & Ausgabe des Hauptprozesses, \R{-f} mitlaufend \\
\R{top <cnt>} & Prozesse im Container \\
\R{cp <quelle> <ziel>} & Dateien Host $\leftrightarrow$ Container \\
\R{commit <cnt> <tag>} & Container einfrieren $\to$ neues Image \\
\R{rm <cnt>}, \R{container prune} & Container löschen / alle gestoppten \\
\addlinespace
\multicolumn{2}{@{}l}{\textbf{Volumes}}\\
\R{volume ls}, \R{volume inspect} & anonyme und benannte Volumes \\
\R{volume prune} & \textbf{löscht Daten} -- siehe Falle \ref{rez:volume} \\
\end{ztab}

\subsection{Dockerfile-Instruktionen}
\begin{ztab}{@{}lXc@{}}
 & Wirkung & Layer \\ \midrule
\R{FROM} & Parent-Image, muss zuerst stehen; mehrfach $=$ mehrstufig & -- \\
\R{COPY} & Datei/Verzeichnis aus dem Kontext ins Image & ja \\
\R{ADD} & wie \R{COPY}, packt Archive aus & ja \\
\R{RUN} & Kommando im entstehenden Image ausführen & ja\,$^*$ \\
\R{WORKDIR} & Arbeitsverzeichnis, wird angelegt & ja \\
\R{ENV} & Umgebungsvariable & -- \\
\R{CMD} & Standard-Kommando \emph{oder} Argumente für ENTRYPOINT & -- \\
\R{ENTRYPOINT} & festes Kommando, Argumente werden angehängt & -- \\
\end{ztab}
$^*$ nur wenn es das Dateisystem ändert: \R{RUN echo "x" > f} schon,
\R{RUN echo "x"} nicht.

\subsection{Drei Rechnungen}
\begin{bsp}[Was an der Prüfung gerechnet wird]
\textbf{(a) Platzbedarf.} Ein Image von 118 MB, drei Container darauf,
die je 2 MB Logfiles schreiben: auf der Platte liegen
$118 + 3\cdot 2 = 124$ MB, nicht $3\cdot 118$. \R{docker ps -a -{}-size}
zeigt pro Container nur die 2 MB (\R{SIZE}) und daneben die Image-Grösse
als \R{virtual}.\\[2pt]
\textbf{(b) Layer zählen.} Jede Zeile, die das Dateisystem ändert, zählt --
einschliesslich des Parent-Images. \R{FROM debian} (1) $+$ \R{RUN apt-get
update} (2) $+$ \R{RUN apt-get install} (3) $=$ 3 Layer; die beiden
\R{RUN} mit \R{\&\&} verbunden $=$ 2 Layer \emph{und} kleiner.\\[2pt]
\textbf{(c) Tag ausschreiben.} \R{nginx} meint vollständig
\R{index.docker.io/library/nginx:latest} -- Repository, Source
(\R{library} $=$ offiziell), Name, Tag.
\end{bsp}
